\subsection{Случайная величина}

Случайная величина является ключевым понятием теории вероятностей.\\
Случайная величина -- это некоторое правило, которое случайному исходу эксперимента сопоставляет число. При этом важно, чтобы для любого множества значений этого числа можно было вычислить вероятность того, что случайная величина попадёт в это множество. Если более строго:

\begin{definition}[Случайная величина]
Случайной величиной в вероятностном пространстве $(\Omega, \mathfrak{S}, \mathbb{P})$ называется измеримая функция $\xi: \Omega \to \mathbb{R}$.
\end{definition}


\textbf{Примеры:}
\begin{itemize}
    \item Бросаем монету: $\Omega = \{\text{орёл}, \text{решка}\}$. Случайная величина $X$: орёл $\to 1$, решка $\to 0$.
    \item Бросаем кубик: $X$ -- выпавшее число очков.
    \item Измеряем рост случайного человека: $X$ -- рост в сантиметрах.
\end{itemize}

Заметим, что множество возможных значений случайной величины тоже
может являться пространством элементарных исходов\\

Случайные величины бывают дискретными, непрерывными и смешанными.

\subsubsection{Распределение случайной величины}

\begin{definition}[Распределение случайной величины]
Распределением случайной величины $X$ называется вероятностная мера $\mu_X$ на числовой прямой $\mathbb{R}$, задаваемая формулой
\[
\mu_X(B) = \mathbb{P}(X \in B) = \mathbb{P}(\{\omega \in \Omega : X(\omega) \in B\})
\]
для любого борелевского множества $B \subseteq \mathbb{R}$.
\end{definition}

Иными словами, распределение -- это правило, которое каждому множеству значений $B$ ставит в соответствие вероятность того, что $X$ примет значение из этого множества.\\

$\xi^{-1}(B)$ Это прообраз множества $B$ при отображении $\xi$. То есть это множество всех элементарных исходов $\omega \in \Omega$, для которых случайная величина $\xi$ принимает значения из множества $B$. Формально:
\[
\xi^{-1}(B) = \{\omega \in \Omega \mid \xi(\omega) \in B\}.
\]

Распределение $\mathbb{P}_{\xi}$ позволяет нам ``забыть'' про исходное пространство $\Omega$ и работать только с числовой прямой. Нам не важно, что именно происходит в $\Omega$ -- важно только, какие значения принимает $\xi$ и с какими вероятностями. Это переход от абстрактного вероятностного пространства к конкретным числам, с которыми удобно работать.\\

Следует упомянуть также следующее определение:
\begin{definition}[Порождённая случайной величиной $\sigma$-алгебра]
Порождённой случайной величиной $\xi$ $\sigma$-алгеброй называется $\sigma$-алгебра $\mathfrak{S}$, состоящая из множеств $\xi^{-1}(B)$, где $B \in \mathfrak{B}(\mathbb{R})$.
\end{definition}

\subsubsection{Функция распределения случайной величины}



Основная характеристика случайной величины -- функция распределения:

\begin{definition}[Функция распределения]
Функцией распределения случайной величины $\xi$ называется функция $F_{\xi}: \mathbb{R} \to \mathbb{R}$ такая, что
\[
F_{\xi}(x) = \mathbb{P}_{\xi}((-\infty; x)).
\]
\end{definition}

\begin{notice}
Выделим основные свойства функции распределения:
\begin{enumerate}
    \item $0 \leqslant F_{\xi}(x) \leqslant 1$, причём $\lim_{x \to +\infty} F_{\xi}(x) = 1$ и $\lim_{x \to -\infty} F_{\xi}(x) = 0$;
    \item $F_{\xi}$ -- возрастающая функция;
    \item $F_{\xi}$ -- непрерывная слева функция, т.е. $\forall c \in \mathbb{R}$
    \[
    \lim_{x \to c-0} F_{\xi}(x) = F_{\xi}(c);
    \]
    \item $\mathbb{P}(\{x_1 \leqslant \xi \leqslant x_2\}) = \lim_{x \to x_2+0} F_{\xi}(x) - F_{\xi}(x_1)$.
\end{enumerate}
\end{notice}

Докажем свойства 3 и 4, поскольку они не столь очевидны.\\

Для этого сначала следует вспомнить аксиому непрерывности
\begin{definition}[Аксиома непрерывности вероятностной меры]
Вероятностная мера $\mathbb{P}$ обладает свойством непрерывности, которое формулируется в двух эквивалентных вариантах:

\begin{enumerate}
    \item \textbf{Непрерывность сверху (для убывающей последовательности):} Если $A_1 \supseteq A_2 \supseteq A_3 \supseteq \ldots$ -- убывающая последовательность событий, то
    \[
    \mathbb{P}\!\left(\bigcap_{n=1}^{\infty} A_n\right) = \lim_{n \to \infty} \mathbb{P}(A_n).
    \]
    В частности, если $\bigcap_{n=1}^{\infty} A_n = \varnothing$, то $\lim_{n \to \infty} \mathbb{P}(A_n) = 0$.

    \item \textbf{Непрерывность снизу (для возрастающей последовательности):} Если $A_1 \subseteq A_2 \subseteq A_3 \subseteq \ldots$ -- возрастающая последовательность событий, то
    \[
    \mathbb{P}\!\left(\bigcup_{n=1}^{\infty} A_n\right) = \lim_{n \to \infty} \mathbb{P}(A_n).
    \]
\end{enumerate}
\end{definition}

\begin{notice}
Аксиома непрерывности является следствием счётной аддитивности вероятностной меры и фактически эквивалентна ей.
\end{notice}

\textit{3.} Используем определение предела по Гейне. Рассмотрим произвольную последовательность $x_1 < x_2 < \ldots < x_n < \ldots < c$ такую, что $\lim_{n \to \infty} x_n = c$. Сопоставим с этой последовательностью следующий набор событий: $A_1, A_2, \ldots$, такой, что $A_n = \{x_n \leqslant \xi < c\}$. Заметим тогда, что $A_1 \supseteq A_2 \supseteq \ldots$ и $\bigcap_{n=1}^{\infty} A_i = \varnothing$. По аксиоме непрерывности имеем
\[
\lim_{n \to \infty} \mathbb{P}(A_n) = 0.
\]
Тогда по Гейне:
\begin{align*}
\lim_{x \to c-0} F_{\xi}(x) &= \lim_{n \to \infty} F_{\xi}(x_n) = \lim_{n \to \infty} \mathbb{P}(\{\xi < x_n\}) = \\
&= \lim_{n \to \infty} \left(\mathbb{P}(\{\xi < c\}) - \mathbb{P}(A_n)\right) = F_{\xi}(c) - \lim_{n \to \infty} \mathbb{P}(A_n) = F_{\xi}(c). \hfill \blacksquare
\end{align*}

\textit{4.}Рассмотрим произвольную последовательность $y_1 > y_2 > \ldots > y_n > \ldots > x_2$ такую, что $\lim_{n \to \infty} y_n = x_2$ (т.е. $y_n \to x_2+0$). Сопоставим с этой последовательностью следующий набор событий: $B_1, B_2, \ldots$, такой, что $B_n = \{x_1 \leqslant \xi < y_n\}$. Заметим тогда, что $B_1 \supseteq B_2 \supseteq \ldots$ и
\[
\bigcap_{n=1}^{\infty} B_n = \{x_1 \leqslant \xi \leqslant x_2\}.
\]
Действительно, если $\omega$ принадлежит всем $B_n$, то $x_1 \leqslant \xi(\omega) < y_n$ для всех $n$, и в пределе при $n \to \infty$ получаем $x_1 \leqslant \xi(\omega) \leqslant x_2$. Обратно, если $x_1 \leqslant \xi(\omega) \leqslant x_2 < y_n$, то $\omega \in B_n$ для всех $n$.

По аксиоме непрерывности имеем
\[
\mathbb{P}\!\left(\bigcap_{n=1}^{\infty} B_n\right) = \lim_{n \to \infty} \mathbb{P}(B_n).
\]
Тогда
\begin{align*}
\mathbb{P}(\{x_1 \leqslant \xi \leqslant x_2\}) &= \lim_{n \to \infty} \mathbb{P}(B_n) = \lim_{n \to \infty} \mathbb{P}(\{x_1 \leqslant \xi < y_n\}) = \\
&= \lim_{n \to \infty} \left(\mathbb{P}(\{\xi < y_n\}) - \mathbb{P}(\{\xi < x_1\})\right) = \\
&= \lim_{n \to \infty} F_{\xi}(y_n) - F_{\xi}(x_1) = \lim_{x \to x_2+0} F_{\xi}(x) - F_{\xi}(x_1). \hfill \blacksquare
\end{align*}

\subsubsection{Важные примечания}

Иногда мы будем использовать понятие независимости случайных величин, однако подробнее его рассмотрим чуть позже

\begin{definition} Случайные величины \(\xi_1\) и \(\xi_2\) называются \textbf{независимыми}, если для любых борелевских множеств \(B_1\) и \(B_2\) верно:
\[
\mathbb{P}(\{\xi_1 \in B_1\}) \mathbb{P}(\{\xi_2 \in B_2\}) = \mathbb{P}(\{\xi_1 \in B_1\} \wedge \{\xi_2 \in B_2\})
\]
\end{definition}

Попарная и совокупная независимость случайных величин определяется аналогично событиям.\\

\begin{definition}[Попарно независимые случайные величины]
Случайные величины $\{\xi_i\}_{i \in I}$ называются \textbf{попарно независимыми}, если любая пара случайных величин из этого набора независима, т.е. для любых $i, j \in I$ ($i \neq j$) и любых борелевских множеств $B_1, B_2 \in \mathfrak{B}(\mathbb{R})$ выполняется:
\[
\mathbb{P}(\xi_i \in B_1, \, \xi_j \in B_2) = \mathbb{P}(\xi_i \in B_1) \cdot \mathbb{P}(\xi_j \in B_2).
\]
\end{definition}

\begin{definition}[Независимые в совокупности случайные величины]
Случайные величины $\{\xi_i\}_{i \in I}$ называются \textbf{независимыми в совокупности} (или просто независимыми), если для любого конечного подмножества индексов $\{i_1, i_2, \ldots, i_n\} \subseteq I$ и любых борелевских множеств $B_1, B_2, \ldots, B_n \in \mathfrak{B}(\mathbb{R})$ выполняется:
\[
\mathbb{P}(\xi_{i_1} \in B_1, \, \xi_{i_2} \in B_2, \, \ldots, \, \xi_{i_n} \in B_n) = \prod_{k=1}^{n} \mathbb{P}(\xi_{i_k} \in B_{i_k}).
\]
\end{definition}

\medskip

Для непрерывных случайных величин можно ввести кроме функции распределения еще и плотность распределения случайной величины, но об этом мы поговорим чуть позже.